% !TeX program = xelatex
\documentclass[UTF8,fontset=fandol,12pt]{ctexart}
\usepackage[a4paper,margin=22mm,headheight=16pt]{geometry}
\usepackage{booktabs,tabularx,array,xcolor,fancyhdr,hyperref}
\definecolor{accent}{HTML}{28649A}
\hypersetup{colorlinks=true,linkcolor=accent,urlcolor=accent,pdftitle={第8队研讨汇报：补充说明与答辩准备},pdfauthor={NLP 第8队}}
\pagestyle{fancy}\fancyhf{}
\fancyhead[L]{\small 第 8 队 · 第一次课设研讨}
\fancyhead[R]{\small 补充说明与答辩准备}
\fancyfoot[C]{\thepage}
\renewcommand{\headrulewidth}{0.3pt}
\setlength{\parindent}{0pt}
\setlength{\parskip}{5pt}
\linespread{1.25}
\newcommand{\question}[1]{\par\vspace{2mm}{\color{accent}\bfseries #1}\par}
\begin{document}
{\LARGE\bfseries 补充说明与答辩准备}\par
2026 年 10 月 22 日研讨\quad 内容更新至 2026 年 10 月 21 日

本文件供答辩前准备和现场查阅，与 12 页幻灯片及 8 分钟讲稿分开。每问可先用两三句话回答，再按老师追问补充细节。

\question{1. 六篇论文与当前项目分别是什么关系？}
六篇论文构成阅读与比较的基础，当前实现以 StructMem 为主线，没有完整复现六套系统。\par\vspace{3mm}
\begin{tabularx}{\linewidth}{@{}p{48mm}X@{}}
\toprule 论文 & 在本项目中的作用\\\midrule
StructMem & 事件绑定与跨事件综合的主方法。\\
Hindsight、Semantic XPath & 为记忆检索与访问方式提供比较。\\
MetaMem & 支撑证据利用与错误分析。\\
HyMEM、Text2Mem & 为记忆更新、操作规范与结果校验提供讨论依据。\\\bottomrule
\end{tabularx}\par\vspace{3mm}

\question{2. 是否严格复现了 StructMem 原论文？}
属于有明确配置差异的方法复现。回答模型使用 DeepSeek-Flash，嵌入使用本地 MiniLM；关闭压缩与话题切分，部分绑定时间由顺序构造，不当作真实消息发送时刻。数据规模和模型条件与论文不同，因此不直接比较本组和论文的总分。

\question{3. S0 是论文中的 flat 消融吗？}
S0 检索原始对话块；论文的 flat 消融使用平铺事实，二者不能混称。当前 S0/S2 的记忆表示、输入长度和方法指令都不同，分差描述的是两套系统。S1 将与 S2 共享事件库和方法指令，只关闭跨事件综合，用于检查综合模块的作用，同时报告输入长度变化。

\question{4. 当前创新性体现在哪里？}
已完成的工作是方法适配、固定题集上的系统对照和证据分析。现有结果提出了具体问题：结构化记忆的回答变化，需要付出多少上下文与构建成本？后续先补综合模块消融，再将检索事件数从 60 调为 30、综合条数保持 5，比较质量与输入量。优化尚未运行，暂不声称提出新算法或已经获得优化收益。

\question{5. 后续怎样评价泛化？}
目前只使用固定的 3 个开发会话、30 道题。开发集上确定配置后，计划从另外 7 个留出会话分层抽取约 100 题评价。保留预定题单、失败和负结果，不根据分数临时换题。

\newpage
{\Large\bfseries 结果、评分与证据边界}\par

\question{6. 63.3\% 能否证明结构化记忆更好？}
应用已有证据裁决后，S0 为 15/30，S2 为 19/30。时间题多答对 2 题，开放域和单跳各多 1 题，多跳相同。这只描述固定开发集上的现象。尚缺 S1 消融和独立留出评价，不能据此作稳定提升或单个模块的因果结论。

\question{7. 为什么 S0 同时出现 16/30 和 15/30？}
16/30 是原始自动裁判结果。一条日期回答同时给出“6 月 24 日”和“此前的星期五”，而记录当天的原文说的是 last Friday，指向 6 月 17 日。已有证据裁决将该题从 yes 改为 no，得到 15/30。原始判定与后续裁决分别保留，S2 仍为 19/30。

\question{8. 评分结果怎样核对？}
已核对 60 个输入指纹，将评分与冻结问题、参考答案和实际预测逐条对应，并离线复算总体与分类分数。原始自动裁判与案例证据裁决分别保留。统计一致性和答案是否符合原文证据分别检查，同源模型回答与裁判也可能存在偏差。

\question{9. 91 条综合记忆已经核查到什么程度？}
已核对综合汇总、事件来源与实际回答输入，并检查相关数值和来源绑定。没有独立回放全部 91 次真实综合调用的原始响应，因此不能把来源及数值一致性扩大为完整的语义验收。

\question{10. 4.42 倍是否等于项目费用增加 4.42 倍？}
4.42 倍是回答阶段的输入 token 比值：S0 为 132,531，S2 为 586,282。它不包含记忆抽取、综合与评分的全部开销，也不是项目总费用倍数。费用估算未核对账户账单，不能称为实际总账单。

\question{11. 为什么不删掉参考证据不足的题？}
6 个核查案例中有 2 题的参考证据不足，涉及游戏平台和海滩住址。继续保留冻结题目、参考答案及分母，并单列限制，避免看见结果后选择性删题。不能把这些题单独用来证明系统检索失败。

\question{12. 目前大致贡献如何理解？}
表中给出当前阶段的大致贡献估计。A 负责数据环境，B 负责普通检索，C 负责事件记忆，D 负责综合与 S2，E 负责评分分析，F 负责统一回答与联调。

\newpage
{\Large\bfseries 小组协作与系统性分析}\par

\question{13. 小组主要通过哪些方式协作？}
主要通过 GitHub、微信群和线下讨论。GitHub 管理任务、代码与文档，并记录交接、审核和修改；微信群用于日常沟通、进度同步和及时协调；线下讨论集中交流论文联系、实验方案、结果与汇报内容。

\question{14. 怎样把六个人的工作衔接成一个整体？}
六人围绕同一批输入和逐题输出交接。A 冻结数据划分和题单，B/C 提供不同记忆表示，D 接收 C 的产物完成综合与 S2，F 提供统一回答入口并联调，E 接收真实回答评分，F 再复算评分。共同题单、输入输出约定和来源记录连接各模块，论文比较则为方法选择和结果分析提供共同框架。

\question{15. 六篇论文的融合具体落实在哪里？}
系统性体现在共同问题、机制联系和设计取舍三个层面。StructMem 与 Semantic XPath 分别帮助检查存储时的关系和读取时的范围；Hindsight 与 MetaMem 帮助区分时间表示、检索和证据利用问题；HyMEM 与 Text2Mem 分别回答更新策略及操作结果如何验证。第 8 页用同一个真实时间案例连接三种视角，不靠逐篇罗列结论。

当前落实的是 StructMem 主实现、来源与会话检查、时间案例分析以及独立评分复算。后续先补综合消融和预算对照；动态更新与反馈规则暂不实施。这是跨论文提炼分析框架和实验问题，不是宣称六套算法已经集成，或把本组验证脚本称为 Text2Mem 的实现。

\vspace{2mm}
依据：六篇论文比较、现有系统性综述、项目方案、协作约定与评分复算。论文页码和图表编号见 README。
\end{document}
